\section{Psychology, Discovery, and Qualification}

\subsection{Sales discovery should reconstruct an actual episode}

A strong opening question is not ``What challenges keep you up at night?'' but ``Tell me about the last time this problem caused a concrete failure or costly workaround.'' Ask what happened, who noticed, what the team did, what it cost, why the current solution was accepted, and what changed now. The buyer's chronology distinguishes an acute, operational problem from a vague, general aspiration.

For an engineer, the episode could be a retrieval-augmented generation (RAG) system returning unreliable or hallucinatory answers to production customers. For a college, it might be a failed elective rollout or employer complaints regarding capstone project quality. For a working professional, it may be an unsatisfactory previous online course that failed to deliver code-level review. The episode supplies observable, empirical evidence: deadlines, projects, currency amounts, hours, missed opportunities, and actual workarounds. Avoid suggesting the desired answer before the narrative is fully told.

\subsection{Jobs-to-be-done focuses the conversation on progress}

People buy tools and educational services in order to make progress in a specific, bounded situation. The software engineer does not merely buy ``AI videos''; they seek to become capable of defending and deploying a production-grade architecture under scrutiny from their engineering director. The enterprise buyer does not merely buy ``corporate training''; they want fewer failed internal AI experiments, faster evaluation cycles, and reduced operational dependence on a single scarce specialist.

A job description must be formulated with sufficient precision to reveal a testable success criterion:
\begin{itemize}
\item \textbf{Vague aspiration:} ``Become future-ready in artificial intelligence.'' (Untestable, impossible to verify during procurement or delivery.)
\item \textbf{Precise economic job:} ``Independently design, benchmark, and evaluate a retrieval system under specified latency and hallucination constraints within 8 weeks.'' (Verifiable, testable, and contractually defensible.)
\end{itemize}
The salesperson must also verify whether the customer values that specific job enough to protect scarce calendar hours and capital. A technically impressive puzzle may not represent a commercially urgent priority.

\subsection{Segment and persona affect both value and process}

The identical technical service is evaluated through radically different lenses by end users, technical evaluators, economic buyers, and organizational sponsors. In Blume Ventures' \textit{GTM Field Notes: B2B Edition} (2025), Ravi Kaushik Rajagopal emphasizes the discipline of mapping persona, problem, and business impact:
\begin{itemize}
\item \textbf{The Engineering Lead:} Evaluates technical validity, architectural reliability, code quality, and reviewer rigor.
\item \textbf{The Finance Owner:} Evaluates total monetary outlay, cash flow timing, and measurable return on investment or cost displacement.
\item \textbf{The Information Security Reviewer:} Evaluates data exposure, 3rd-party model access, intellectual property leakage, and compliance standards.
\item \textbf{The Executive Sponsor:} Evaluates enterprise execution risk, time-to-market acceleration, and strategic talent retention.
\end{itemize}
A founder must never deliver mutually contradictory promises to these distinct stakeholders. Instead, translate the single underlying operational mechanism into its relevant consequences for each role. Technical validity supports the economic case; clear scope and governance documentation support procurement; delivery reliability reassures the executive sponsor. If the founder cannot identify who possesses the authority to sign an invoice, the deal remains in early discovery regardless of user excitement.

\subsection{Qualification is not a rigid acronym}

Enterprise frameworks such as BANT (Budget, Authority, Need, Timing), MEDDIC (Metrics, Economic Buyer, Decision Criteria, Decision Process, Identify Pain, Champion), and their modern derivatives help prevent critical governance questions from being overlooked. However, they are not substitutes for situational judgment. Some early adopters have urgent operational pain but no preallocated line-item budget; they can create 1 if the business case is compelling. Other buyers possess ample budget and authority but lack acute pain, resulting in endless polite conversations that never close.

For a lean founder, 5 foundational questions suffice:
\begin{enumerate}
\item What is the costly, unaddressed problem?
\item Why is it imperative to act now rather than 6 months from now?
\item Who possesses the formal authority to authorize expenditure and sign contracts?
\item Which alternative will be chosen if no transaction occurs?
\item What specific evidence and contractual conditions would make commitment rational?
\end{enumerate}
Record the answers as empirical facts with attributed sources. A label such as ``qualified'' in a spreadsheet without documented answers is merely a salesperson's optimistic opinion.

\subsection{A decision trigger is stronger than a demographic category}

A prospect who matches an ideal demographic profile (for example, a 28-year-old software engineer in Bangalore) may have zero motivation to purchase education today. The identical engineer becomes urgently motivated upon being assigned to lead a production machine learning deployment or after receiving critical feedback during an annual technical review. The trigger creates operational relevance and immediate commercial urgency.

In institutional sales, triggers are equally decisive:
\begin{itemize}
\item A college revisits its curriculum when academic committee windows open, annual budgets are allocated, accreditation visits are scheduled, or placement statistics decline.
\item An enterprise funds workflow training when a major software release is blocked by model evaluation failures.
\end{itemize}
Sales qualification must identify authentic decision triggers rather than manufacture artificial urgency. Observing that ``AI is transforming the global economy'' is a broad macro observation; discovering that ``your engineering team faces a production reliability audit in 8 weeks'' provides an actionable commercial foundation.

\subsection{The cost of inaction requires careful estimation}

Consider an enterprise engineering organization where senior developers lose 15 hours each week conducting repetitive, manual data verification and hallucination debugging across internal prototypes. At an internally estimated blended labor cost of ₹1,500 per hour, the gross annualized effort expended on this workaround is calculated as:
\begin{equation}
\begin{aligned}
\text{Gross Burden} &= 15 \text{ hrs/wk} \times ₹1,500/\text{hr} \times 52 \text{ wks} \\
&\approx ₹11,70,000 \quad \text{(₹11.7 lakh)}.
\end{aligned}
\end{equation}
However, the consultative seller must exercise analytical caution: not every recovered hour converts directly into cash savings, and no educational or software intervention eliminates 100\% of operational friction.

The seller must explicitly decompose gross wasted effort into realistic economic categories:
\begin{itemize}
\item \textbf{Cash Savings:} Direct reduction in external contractor fees or compute infrastructure waste.
\item \textbf{Capacity Release:} Hours redirected toward revenue-generating product features (realized only if management reallocates time effectively).
\item \textbf{Risk Mitigation:} Downside avoidance from customer churn or SLA penalties caused by bad model outputs.
\end{itemize}
In consultative selling, use the buyer's own conservative parameters, state all assumptions explicitly, and invite scrutiny. A vendor-generated ``ROI calculator'' built on aggressive multipliers instantly destroys credibility. Present a realistic range: baseline burden, addressable fraction (typically 30\% to 60\%), implementation overhead, and downside risk.

\subsection{Listening is an economic tool}

A founder who dominates 80\% of a discovery conversation learns virtually nothing about the buyer's real alternatives or internal constraints. High-yield discovery requires creating deliberate conversational space for the prospect to articulate facts, anxieties, and organizational dynamics. Ask 1 question at a time. Request concrete historical examples. Reflect the customer's problem back accurately and pause to allow correction.

When an enterprise manager says, ``We need an intensive AI upskilling program,'' do not launch into a 20-minute pitch about curriculum modules. Ask:
\begin{quote}
\textit{``Which specific production failure is this program intended to prevent, and why has your internal senior engineering staff been unable to resolve it using existing documentation and code reviews?''}
\end{quote}
The objective of disciplined listening is not to adhere to an arbitrary conversational ratio, but to maximize the rate of verifiable information extraction per minute of founder time.
